% !TEX root = ../../main.tex

\section{功能性需求分析}

\subsection{普通用户功能需求}

普通用户是 HearBridge 移动端的主要使用者，其核心目标是通过移动端完成手语资源学习、实时训练和句子视频识别。
普通用户进入应用后，若尚未登录，需要先完成登录或注册；
若已经登录，则可直接进入应用主界面。
系统主界面包括首页、训练和我的等主要功能入口。
其中，首页主要承担快捷入口和常用内容展示作用，训练页面主要用于浏览或搜索手语资源，我的页面主要用于个人信息和账号相关操作。

普通用户的功能需求主要包括以下几个方面：

（1）用户认证需求。
普通用户应能够完成注册、登录和退出登录等操作。
系统需要在用户启动应用时判断其登录状态，未登录用户不能直接进入主界面，必须先完成登录或注册。
用户退出登录后，系统应清除当前登录状态，并要求用户重新认证后才能继续使用需要登录态的功能。

（2）手语资源浏览需求。
用户应能够查看系统提供的手语分类和手语资源列表，并能够根据训练需求选择具体手语资源。
手语资源应包含名称、分类、封面或相关展示信息，使用户能够在进入训练前了解资源内容。

（3）动作示范查看需求。
用户选择某一手语资源后，应能够进入对应的手势训练页面，并查看该手语动作的示范内容。
动作示范用于帮助用户理解标准动作，为后续模仿练习和实时识别提供参考。

（4）实时手势训练需求。
用户应能够在训练页面开始实时练习，通过摄像头对准自身手势动作，并查看系统返回的实时识别结果。
识别结果应以用户容易理解的形式展示，例如识别标签、置信度或结果提示，使用户能够及时了解当前动作是否接近目标动作。

（5）句子视频识别需求。
用户应能够从首页入口进入句子视频识别页面，选择本地手语视频并提交识别。
识别完成后，系统应向用户展示原始识别结果和语义修正结果，帮助用户理解视频中手语动作对应的表达内容。

（6）个人信息管理需求。
用户应能够在个人中心查看和修改个人信息，并根据需要修改密码或退出登录。
个人信息管理功能应保持操作简洁，避免影响用户使用训练和识别功能的主流程。

\WhuImageFigure[0.92\textwidth][0.46\textheight]{figures/chapter03/user-usecase.pdf}{普通用户用例图}{fig:user-usecase}

如图~\ref{fig:user-usecase} 所示，普通用户可以完成登录、浏览手语资源、查看资源详情、观看动作示范、进行实时手势识别、上传句子视频进行识别、查看识别结果、查看语义修正结果、修改个人信息和退出登录等操作。
这些用例共同构成普通用户在移动端的主要功能范围。

\subsection{管理员功能需求}

系统管理员是 HearBridge 后台管理端的主要使用者，其核心目标是维护移动端训练和识别功能所依赖的资源、样本和模型。
与普通用户关注训练体验不同，管理员更关注数据维护、模型训练、模型版本发布和系统运行状态查看等管理操作。

管理员的功能需求主要包括以下几个方面：

（1）管理员登录需求。
管理员进入后台管理端前，需要完成身份认证。
系统应限制未登录用户访问管理页面，避免未授权人员操作训练资源、样本数据和模型版本。

（2）手势分类管理需求。
管理员应能够维护手势资源所属的分类信息，包括分类名称、分类编码、排序和启用状态等内容。
分类管理用于组织手语资源，便于移动端用户按照分类查找训练内容。

（3）手势资源管理需求。
管理员应能够维护具体手势资源，包括资源编码、中文名称、所属分类、封面资源、数字人动作资源或相关示范资源地址等信息。
系统应支持资源的新增、修改、查询和状态维护，使移动端能够持续获取可训练的手语内容。

（4）样本管理需求。
管理员应能够查看和管理手势识别训练过程中产生或导入的样本数据，包括样本所属标签、采集状态、质量信息和可用状态等。
样本管理用于为后续特征转换、模型训练和错误数据清理提供依据。

（5）模型训练管理需求。
管理员应能够发起或查看模型训练任务，了解训练过程和训练结果。
模型训练管理用于支持识别模型持续迭代，使系统能够根据新增样本和扩展词表更新识别能力。

（6）模型版本管理需求。
管理员应能够维护模型版本信息，包括版本名称、模型文件、支持动作数量、训练结果和发布状态等。
系统需要区分不同模型版本，便于后续回溯、切换和发布。

（7）发布模型需求。
管理员应能够将某一模型版本发布为当前识别服务使用的模型。
发布后，移动端实时识别和句子视频识别功能应能够使用当前发布模型或相关配置进行识别。

（8）系统概览需求。
管理员应能够查看系统概览信息，例如资源数量、样本数量、模型版本状态或其他运行概况，以便快速了解系统当前维护状态。

\WhuImageFigure[0.92\textwidth][0.52\textheight]{figures/chapter03/admin-usecase.pdf}{管理员用例图}{fig:admin-usecase}

如图~\ref{fig:admin-usecase} 所示，系统管理员可以完成管理员登录、手势分类管理、手势资源管理、样本管理、模型训练管理、模型版本管理、发布模型、查看系统概览和退出登录等操作。
这些功能共同支撑移动端训练内容维护和识别模型迭代。
